\clearpage
\appendix

\setcounter{section}{0}
\setcounter{table}{0}
\setcounter{figure}{0}
\renewcommand{\thesection}{S\arabic{section}}   
\renewcommand{\thetable}{S\arabic{table}}   
\renewcommand{\thefigure}{S\arabic{figure}}

In this supplementary material, we present additional technical details and more experimental results that could not be included in the main manuscript due to the lack of pages.
The contents are summarized below:
\begin{itemize}
\vspace{2.5mm}
% \item \ref{sec:visual_video}. Visualization in video format
\item \ref{sec:suppl_text_align}. Evaluation on semantic alignment
\item \ref{sec:suppl_hdm}. Comparison with HDM
\item \ref{sec:suppl_contact}. Impact of contact estimation accuracy
\item \ref{sec:suppl_gs_optim}. Impact of Gaussian attributes optimization
\item \ref{sec:suppl_g2m}. Impact of Gaussians-to-mesh conversion
\item \ref{sec:suppl_captioning}. Details of text captioning
\item \ref{sec:suppl_implementation}. Implementation details
\item \ref{sec:suppl_limitation}. Limitations and future work
\item \ref{sec:suppl_more_results}. More qualitative results
\end{itemize}


\subsection{Details of video decomposition}
\label{sec:supp_decomposition}

This section expands \cref{sec:3.2_decomposition}: the exact signal used to
propose sub-action boundaries, the motion evidence handed to the captioner, the
prompts of every language-model call, and the post-processing that fixes the
final sub-action lengths.

\noindent\textbf{Boundaries from the object trajectory.}
The relative translation $\mathbf{p}_t$ of \cref{eq:relative_translation} is
built from the initial human and object motions alone, so it is available before
any generation or optimization takes place.  We smooth it with a moving average
over $w=5$ frames, replicating the end frames, which removes single-frame
tracker jitter without displacing a real regime change by more than two frames.
The recursion of \cref{eq:rdp} accepts a boundary when its residual exceeds
$\epsilon=0.1$\,m, and never proposes one within $g=15$ frames of either end of
the segment it splits; a segment shorter than $2g$ is therefore left intact.
Because the residual is measured against the straight line through the two
endpoints, a segment in which the object is carried at a constant relative
velocity is never split, however long it is, while a segment that reverses or
stops is split at the reversal.  When either the human or the object track is
missing, we fall back to uniform windows of $81$ frames, which is the clip
length that the generation stage of \cref{sec:3.3_generation} consumes.

\noindent\textbf{Body parts for the motion evidence.}
$\mathcal{B}$ contains the seven SMPL-X joints listed in
\cref{tab:supp_body_parts}.  The set is deliberately small: over $278$ segments
with clean tracking, every joint we exclude follows one of these seven with a
Pearson correlation of at least $0.93$ (neck with shoulder $0.98$, elbow with
its own hand $0.95$, the two knees $0.96$), so adding them would repeat lines of
the prompt rather than add information.  The two hands are kept apart because
they are the one genuinely independent pair ($r=0.63$), as one hand acts while
the other rests.  Hips are kept despite following the chest at $r=0.96$, since a
sizable part of our data is furniture that is sat on, and \textit{hips} is the
word that puts the captioner there.

\begin{table}[t]
\centering
\small
\begin{tabular}{lc}
\toprule
Reported part & SMPL-X joint index \\
\midrule
Head        & 15 \\
Chest       & 6 \\
Hips        & 0 \\
Left hand   & 20 \\
Right hand  & 21 \\
Left foot   & 10 \\
Right foot  & 11 \\
\bottomrule
\end{tabular}
\caption{The body parts $\mathcal{B}$ whose distance change to the object is
reported to the captioner, and the SMPL-X joints they are measured at.}
\label{tab:supp_body_parts}
\end{table}

\noindent\textbf{Verbalizing the distance change.}
Each $\Delta_{k,j}$ of \cref{eq:part_proximity} becomes one of three words:
\textit{unchanged} when $|\Delta_{k,j}|<\tau$ with $\tau=10$\,cm,
\textit{moves closer} when $\Delta_{k,j}\le-\tau$, and \textit{moves away}
otherwise.  On the segments above, this band leaves $27\%$ of the readings
\textit{unchanged}, which keeps that verdict meaningful without swallowing real
approach and retreat.  Only the first and the last frame of a segment are
compared, which the prompt states explicitly, so a part that drifts away and
comes back still reads as having moved closer.  The resulting block, which is
substituted verbatim into \cref{fig:supp_prompt_segment}, reads

\begin{quote}\small\begin{verbatim}
Measured change in the object's 3D distance to
each body part, first frame of this segment vs
last:
- Head: moves closer
- Chest: moves closer
- Hips: unchanged
- Left hand: unchanged
- Right hand: moves closer
- Left foot: moves away
- Right foot: moves away
\end{verbatim}
\end{quote}

\noindent\textbf{Withholding the evidence.}
A single frame at which the tracker has lost the object corrupts every distance
measured across the window that contains it, so we test the raw, unsmoothed
object pose for abrupt steps between adjacent frames and emit an empty evidence
block for any segment that contains one.  A step is abrupt when the translation
moves by more than $50$\,cm in one frame, or when the rotation turns by more than
$90^\circ$ in one frame; both frames spanning the step are marked.  The
rotation test is applied only to sequences whose median per-frame turn is below
$5^\circ$, that is, only where the tracker resolves rotation at all: on round or
thin objects it does not, and such sequences show a median turn near $14^\circ$
with a fifth of frames beyond $150^\circ$, which is an unresolved rotation
rather than a spinning object.  Both bars are loose by design, as they exist to
catch tracking failure and not to judge fast motion.  Scaling them to each
sequence's own median step was tried and is worse, because a sequence whose
tracking fails throughout raises its own bar and passes.

\noindent\textbf{Per-segment captioning.}
Every segment produced by the recursion above is captioned independently with
Qwen3-VL-8B-Instruct~\cite{bai2025qwen3}.  We sample up to $5$ frames uniformly
from the segment, resize them so that the longer side is $224$ pixels, annotate
each with its timestamp, and prepend the system line \textit{``You are a precise
video annotator for hand-object interaction segments.''} to the prompt of
\cref{fig:supp_prompt_segment}.  Decoding is greedy with at most $128$ new
tokens.  The prompt forbids the words \textit{and} and \textit{then} so that one
caption names one action; when a caption still contains them, we keep its longest
clause.

\begin{figure*}[t]
\centering
\small
\begin{verbatim}
Analyze these ordered, uniformly sampled frames from one segment of a monocular human-object
interaction video.

The interacted object is {object_name}.
The segment covers {full_start:.2f}s to {full_end:.2f}s in the full video ({duration:.2f} seconds
long). You are given {sample_count} uniformly sampled frames from it.
{motion_hint}

Task:
- Write exactly one short English caption describing the dominant human-object interaction across
  these frames.
- Refer to the object as "{object_name}".
- Prefer a motion or state change over a static appearance description.
- Express exactly one action; do not use the words "and" or "then".
- Do not include timestamp phrases, frame numbers, or generic text like "the interaction continues".

Examples (format and range of detail only; do not copy these actions):
- {"caption": "A person holds the {object_name} against the chest."}
- {"caption": "A person lifts the {object_name} with one hand."}
- {"caption": "A person sets the {object_name} down."}
- {"caption": "A person rotates the {object_name} slowly while gripping it near one end with the
   right hand."}
- {"caption": "A person pushes the {object_name} away from the body across the surface with an open
   palm."}
- {"caption": "A person lowers the {object_name} onto the surface with the right hand, steadying it
   with the left."}

Output Format:
{"caption": str}
\end{verbatim}
\caption{Per-segment captioning prompt.  \texttt{\{motion\_hint\}} is the
measured-motion block described in \cref{sec:supp_decomposition}, and is empty
whenever the tracked pose may not be quoted.  \texttt{\{object\_name\}} is the
object category of the clip.}
\label{fig:supp_prompt_segment}
\end{figure*}

\noindent\textbf{Temporal grounding.}
Each caption is grounded independently against the whole clip with
TimeLens-7B~\cite{zeng2025timelens}, using the query

\begin{quote}\small\begin{verbatim}
Please find the visual event described by the
sentence '{caption}', determining its starting
and ending times. The format should be: 'The
event happens in <start time> - <end time>
seconds'.
\end{verbatim}
\end{quote}

\noindent
The video is sampled at $6$\,fps with a visual-token budget of $3584$ tokens,
decoding is greedy with at most $128$ new tokens, and the returned interval is
clamped to the clip and ordered.  Two consecutive grounded intervals are merged
when their temporal intersection over union reaches $0.9$, which collapses only
near-duplicate groundings and leaves a contained short event as its own segment.
From the merged result we keep the start times only, pin the first one to frame
$0$, and let each sub-action end one frame before the next one starts, so the
final sub-actions form a contiguous partition of the clip with no gap or
overlap.

\noindent\textbf{Joint re-captioning.}
The per-segment captions are written without knowledge of one another, so
adjacent segments can receive the same sentence.  We therefore re-caption the
grounded timeline in a single pass with the prompt of
\cref{fig:supp_prompt_merge}, which sees up to $32$ frames sampled at $2$\,fps
over the full clip together with all preliminary captions and their intervals,
and which is required to give every segment a distinct description.  Decoding is
greedy with at most $1536$ new tokens, since the answer is a JSON array over the
whole clip; a reply that does not parse is retried up to three times with an
explicit reminder of the schema, and a truncated reply falls back to the
segments that were emitted completely.  The number of sub-actions is left to the
model unless a target count is supplied, in which case the prompt requests that
count and the grounding merge above is given the same budget.

\begin{figure*}[t]
\centering
\small
\begin{verbatim}
Analyze this merged monocular human-object interaction video segment.

The interacted object is {object_name}. The full clip is {duration:.2f} seconds long.

Preliminary segment results are included below.
They are useful evidence for short actions that may be hard to see in sparse samples, but their
boundaries can be imperfect.

Preliminary segments:
{source_captions}

Task:
- {subaction_count_instruction}
- Segment the complete clip in chronological order, using the frames and preliminary captions
  together.
- Express exactly one action in each caption; do not use the words "and" or "then".
- Every segment must have a distinct caption describing its own visible action or state change;
  never repeat, paraphrase, or reuse the same action description across segments.
- If adjacent segments appear similar, identify the most specific visible difference between them
  instead of assigning the same caption.

Examples:
- {"subactions": [
  {"caption": "A person turns toward the chair.", "start": 0.0, "end": 0.7},
  {"caption": "A person lowers the body toward the chair seat.", "start": 0.7, "end": 1.6},
  {"caption": "A person settles onto the chair with both feet on the floor.", "start": 1.6,
   "end": 2.5}
]}

Output Format:
{"subactions": [{"caption": str, "start": float, "end": float}]}
\end{verbatim}
\caption{Joint re-captioning prompt, abbreviated to one of its three
in-context examples.  \texttt{\{source\_captions\}} lists every preliminary
caption with its grounded interval, and
\texttt{\{subaction\_count\_instruction\}} either requests a given number of
segments or asks for the visually supported number.}
\label{fig:supp_prompt_merge}
\end{figure*}

\noindent\textbf{Sub-action length.}
The boundaries so far follow the interaction, not the generator, so we finally
constrain their spacing.  While any sub-action is shorter than $30$ frames, we
remove the boundary of the shortest one, merging it into the neighbor it was
split from; the boundary at frame $0$ is never removed.  We then split every
sub-action longer than $150$ frames into $\lceil L/150 \rceil$ pieces of equal
length that inherit its caption.  The lower bound keeps a sub-action long enough
for an action to be visible and for grounding to be meaningful, and the upper
bound keeps it within the horizon that the generation stage handles in one pass.
Both bounds are enforced on frames rather than seconds, since the generator
consumes a fixed number of frames.

\noindent\textbf{Contact captioning.}
For each final sub-action, we sample $5$ frames uniformly and query the same
vision-language model with the prompt of \cref{fig:supp_prompt_contact}, which
restricts the answer to a closed vocabulary of $18$ body parts.  We parse the
returned JSON, keep only the names that appear in the vocabulary, and map each
name to a group of SMPL-X body-part vertex segments, \textit{e.g.}, \textit{left
hand} to \texttt{leftHand} and \texttt{leftHandIndex1}, \textit{torso} to
\texttt{spine}, \texttt{spine1}, and \texttt{spine2}, and \textit{right knee} to
\texttt{rightUpLeg} and \texttt{rightLeg}.  The union of these vertices defines
$\mathbf{c}_t\in\{0,1\}^{V}$ over the $V$ SMPL-X vertices, constant within a
sub-action and zero wherever the model reports no visible contact.  Because the
label is derived from the same temporally grounded description that drives
generation, the supervised contact follows the interaction as it changes over
time, which a dense contact map regressed per frame from appearance does not
guarantee.

\begin{figure}[t]
\centering
\small
\begin{verbatim}
Analyze these ordered, uniformly sampled images
from one human-object interaction video chunk.

The interacted object is {object_name}.

Task:
- Identify the body parts that are in direct
  physical contact with the {object_name}
  during this chunk.
- Choose only from this pre-defined list;
  multiple parts are allowed: {parts}
- Focus only on the person whose body center is
  closest to the image center.
- "left" and "right" are the person's anatomical
  sides, not the viewer's sides.
- Include a part only when direct physical
  contact is visible. If no clear contact is
  visible, return an empty list.

Examples:
- {"contact_parts": ["right hand"]}
- {"contact_parts": ["right arm", "right elbow"]}
- {"contact_parts": ["left foot"]}
- {"contact_parts": ["left thigh", "right thigh"]}

Return JSON only, using this exact shape and
replacing "part name" with entries from the
allowed list: {"contact_parts": ["part name"]}
For no visible contact, return:
{"contact_parts": []}
\end{verbatim}
\caption{Contact captioning prompt.  \texttt{\{parts\}} is the closed
vocabulary of $18$ body parts: head, neck, torso, hips, left/right shoulder,
left/right arm, left/right elbow, left/right hand, left/right thigh, left/right
knee, and left/right foot.}
\label{fig:supp_prompt_contact}
\end{figure}

\noindent\textbf{Hyperparameters.}
\cref{tab:supp_decomposition_hparams} collects every constant of this stage.
All of them are shared across datasets; none is tuned per sequence.

\begin{table}[t]
\centering
\small
\begin{tabular}{lc}
\toprule
Smoothing window $w$ & 5 frames \\
Boundary residual threshold $\epsilon$ & 0.1\,m \\
Minimum split gap $g$ & 15 frames \\
Distance tolerance $\tau$ & 10\,cm \\
Abrupt translation step & 50\,cm / frame \\
Abrupt rotation step & $90^\circ$ / frame \\
Rotation-resolved median turn & $5^\circ$ / frame \\
Minimum sub-action length & 30 frames \\
Maximum sub-action length & 150 frames \\
\midrule
Frames per segment caption & 5 @ 224\,px \\
Frames for joint re-captioning & $\le32$ @ 2\,fps \\
Frames per contact query & 5 \\
Grounding frame rate & 6\,fps \\
Grounding token budget & 3584 \\
Grounding merge IoU & 0.9 \\
\bottomrule
\end{tabular}
\caption{Hyperparameters of video decomposition.}
\label{tab:supp_decomposition_hparams}
\end{table}

\subsection{Details of Gaussian re-rendering}
\label{sec:supp_rerendering}

\noindent\textbf{Novel-view rendering.}
The four orbit cameras are placed by rotating the input camera about the
subject's vertical axis by $\{-60^{\circ},-30^{\circ},+30^{\circ},+60^{\circ}\}$,
orbiting the point at the estimated subject depth, and all views are rendered on
a fixed $480\times832$ portrait canvas so that every view and every sequence
share one output geometry.
The Gaussians are posed by the \emph{initial} motion, before any optimization,
so this stage depends only on \cref{sec:3.2_decomposition}.

\noindent\textbf{Restoration.}
We use Wan~2.1~\cite{wan2025wan} at $480$p with its control branch driven by a
DWPose~\cite{yang2023effective} skeleton extracted from the rendering itself.
The restoration is a partial denoising pass at strength $0.7$ with $7$ steps
under an SDE sampler, and the guidance weight over the rendered foreground is
$\lambda=0.05$; the foreground mask is read from the render's own non-background
pixels rather than re-segmented, since we rasterized it and its alpha is already
exact and temporally consistent.
Clips longer than the model's trained length of $81$ frames are denoised in a
single pass with the transformer tiled over time, seeing $21$ latent frames at a
stride of $15$, with the overlaps crossfaded \emph{inside} every denoising step
so that latents mix across seams.

\noindent\textbf{Sub-action chunking.}
Passes are cut on the boundaries $\{(s_k,e_k)\}$ of \cref{sec:3.2_decomposition}.
A segment shorter than $33$ frames is far out of the model's distribution and is
merged into its predecessor, with the absorbed caption appended to the prompt so
that a merged pass still describes both sub-actions.
Adjacent segments share $16$ frames, split evenly on either side of the
boundary, and are joined with a cosine crossfade.
The appearance reference for segment $k$ is the frame at $s_k$ of the real input
video, masked by the human and object silhouettes; a single reference for the
whole clip would keep showing the subject's pose and occlusions at frame $0$
long after it has turned.

\subsection{Details of 4D HOI optimization}
\label{sec:supp_optimization}

We optimize for $25$ steps with Adam ($\beta=(0.9,0.999)$) at learning rates
$10^{-3}$ for the human pose, $10^{-4}$ for the human shape, and $5\times10^{-2}$
for the object rotation and translation.
The motion is parameterized by a cubic B-spline with one coefficient every $9$
frames.
Each step samples $4$ temporal windows of $16$ consecutive frames, drawn so that
window anchors at the sequence edges are visited first and every frame anchors a
window exactly once per epoch; the per-frame losses are then divided by the
number of windows covering that frame.
$\mathcal{L}_{\text{cent}}$ is applied for the first $10$ steps and disabled
afterwards, and the contact threshold is $\tau=1.0$.
Rendering during optimization is performed at half the input resolution.